% The command-line tools: how to run them, their conventions, and one reference entry per
% tool. The option lists are those of the Python implementation (python/src/wfdbx/apps),
% which check_manual.py compares with the usage text of each tool.
\chapter{The command-line tools}
\label{chap:tools}

\wfdbx{} includes the 33 command-line tools of the WFDB Software Package \wfdbversion{}
that read, write and convert records and annotation files. Each of the seven
implementations has all of them, with the options, output, messages and exit status of the
C tools, so that a script written for the WFDB tools runs unchanged with \wfdbx{}. The
exceptions are the deliberate divergences of \cref{chap:behavior}: no built-in signal
parameters (hence the options \opt{-F}, \opt{-G}, \opt{-R} and \opt{-U} added to some
tools), files found only where they are named, annotation files sorted without loss, an
output signal file never written over an input file, write errors reported in the exit
status, and the record line of a header found by \code{getinfo} as the header parser finds
it (the \code{Notes} of \cmd{wfdbdesc}). The original behavior remains available with \opt{-{}-legacy}
(\cref{sec:tools-legacy}).

Every implementation is checked against the C tools by the same conformance suite (1155
cases): the standard output, the standard error, the exit status and the files written are
compared byte for byte. The examples of this chapter were produced by the Python
implementation, run in a copy of the directory \file{tests/data} of the repository, which
holds the sample record \code{100s} (the first minute of record 100 of the MIT-BIH
Arrhythmia Database) and the other test inputs; every other implementation gives the same
output.

\section{Running the tools}
\label{sec:tools-running}

The tools are separate programs in C++ and Rust, and subcommands of a launcher in the other
languages, as in \cref{tab:tools-launchers}. The arguments that follow the tool's name are
those of the C tool: \cmd{wfdbx rdsamp -r 100s -t 10} (Python, Julia, R, MATLAB/Octave and
Kotlin, with the alias \cmd{wfdbx}) and \cmd{rdsamp -r 100s -t 10} (C++ and Rust) do the
same thing. The tool sees its own name as its program name, so its messages begin with
\code{rdsamp:}, as those of the C tool do.

\begin{xltabular}{\linewidth}{@{}P{2.5cm}L@{}}
\caption{How each implementation runs a tool (here \code{rdsamp}).}
\label{tab:tools-launchers}\\
\toprule
{\leavevmode\tabhead Language} & {\leavevmode\tabhead Command} \\
\midrule
\endfirsthead
\toprule
{\leavevmode\tabhead Language} & {\leavevmode\tabhead Command} \\
\midrule
\endhead
\bottomrule
\endfoot
Python & \cmd{wfdbx rdsamp ...}: the command \cmd{wfdbx} is installed with the package
  (\cmd{pip install}); without it, \cmd{python -m wfdbx.cli rdsamp ...} or
  \cmd{python -m wfdbx.apps.rdsamp ...}. In a program, \code{wfdbx.apps.rdsamp.main(argv)}
  runs the tool (\code{argv} is the whole argument list, the program name first) and
  returns its exit status. \\
C++ & \cmd{rdsamp ...}: one program per tool, built in \file{cpp/build/bin/} (or the
  \file{bin} directory of the installation). \\
Rust & \cmd{rdsamp ...}: one program per tool, installed by
  \cmd{cargo install -{}-path rust}. \\
Julia & \cmd{julia -{}-startup-file=no -{}-project=JULIADIR JULIADIR/bin/wfdbx.jl rdsamp ...},
  where \code{JULIADIR} is the directory \file{julia} of the repository; a shell script
  named \cmd{wfdbx} that runs this command makes it \cmd{wfdbx rdsamp ...}. In a session,
  \code{WFDBX.run\_tool("rdsamp", ["-r", "100s"]; dir=...)} returns the exit status and the
  bytes written to the standard output and the standard error, and \code{WFDBX.tools()}
  the names of the tools. \\
R & \cmd{Rscript r/inst/bin/wfdbx.R rdsamp ...} (in an installed package, the script is
  \code{system.file("bin", "wfdbx.R", package = "wfdbx")}); with the alias
  \cmd{alias wfdbx='Rscript .../wfdbx.R'}, \cmd{wfdbx rdsamp ...}. In a session,
  \code{wfdbx\_tool("rdsamp", c("-r", "100s"))} runs the tool and returns its exit status,
  and \code{wfdbx\_tools()} lists the tools. \\
MATLAB/Octave & Octave: \cmd{octave-cli -{}-no-gui -{}-norc -q matlab/bin/wfdbx\_cli.m rdsamp ...}
  (with an alias \cmd{wfdbx}, \cmd{wfdbx rdsamp ...}); MATLAB:
  \cmd{matlab/bin/wfdbx\_matlab.sh rdsamp ...}, which runs \cmd{matlab -batch} (or the
  program named by the environment variable \code{WFDBX\_MATLAB}). In a session,
  \code{status = wfdbx.apps.rdsamp(\{'-r', '100s'\})}; \cmd{help wfdbx.apps.rdsamp} shows
  the usage. \\
Kotlin & \cmd{kotlin/bin/wfdbx rdsamp ...}, a POSIX shell script that runs
  \cmd{java -cp build/libs/wfdbx-all.jar wfdbx.cli.MainKt rdsamp ...} (the jar is built by
  \cmd{./gradlew fatJar}); a link to \file{bin/wfdbx} named after a tool runs that tool. \\
\end{xltabular}

The launchers also take a few arguments of their own, in place of the tool's name:
\opt{-{}-list} prints the names of the tools, one per line (Python, R, MATLAB/Octave and
Kotlin); \opt{-{}-version} prints the version of WFDB-x, as \code{WFDB-x} followed by the
version number (Python, Julia, R, MATLAB/Octave and Kotlin); \opt{-h} or \opt{-{}-help}
prints the launcher's usage (Python, R, MATLAB/Octave and Kotlin), and the Python launcher
also has \opt{-{}-examples} \textit{DIR}, which copies the example notebooks to \textit{DIR}
(by default \file{wfdbx-examples}). Without arguments, or with an unknown tool name, a
launcher prints a message on the standard error and exits with status 1. The programs of
C++ and Rust, like the C tools, have no \opt{-{}-version} option.

\section{Conventions}
\label{sec:tools-conventions}

\subsection{Records and files}
\label{sec:tools-files}

A record is named as in \cref{chap:records}, without the suffix of its header
(\code{100s}, not \code{100s.hea}), and an annotation file by its record and its
annotator (\opt{-r 100s -a atr} is the file \file{100s.atr}). The tools find every file
that they read through the library, as \cref{chap:files} describes: a relative name is
taken relative to the \emph{input directory}, set by the environment variable
\code{WFDBX\_INPUT\_DIR} (by default the current directory), an absolute name, a \code{\textasciitilde}
name or a URL is used as given, and nothing is searched for. A URL (only \code{http://} and
\code{https://} URLs are read) is opened only when network access is turned on, by
\code{WFDBX\_NETWORK=1}. Files written through the library
(headers, signal files, annotation files) are created in the \emph{output directory},
\code{WFDBX\_OUTPUT\_DIR} (by default the current directory). For example, to read record
100 of the MIT-BIH Arrhythmia Database from PhysioNet:
\begin{lstlisting}[style=shell]
export WFDBX_INPUT_DIR=https://physionet.org/files/mitdb/1.0.0/
export WFDBX_NETWORK=1
wfdbx rdsamp -r 100 -t 0:05
\end{lstlisting}

A few files are opened by the tools themselves, not through the library: the input files
named by \opt{-i} in \cmd{wrsamp}, \cmd{edf2mit}, \cmd{wav2mit}, \cmd{a2m} and \cmd{ad2m},
the input files of \cmd{ahaecg2mit}, the script of \cmd{xform -S}, and the signal file
whose size \cmd{wfdbdesc} measures; and the output files that a tool writes itself: the WAV
file of \cmd{mit2wav}, the EDF file of \cmd{mit2edf}, the master header of \cmd{snip -m},
and the \file{.des}, \file{.ecg}, \file{.key} and \file{.txt} files of \cmd{parsescp}.
They follow the same rules as the files of the library (\cref{sec:tool-files}): a relative
name is read from the input directory, or created in the output directory (with its
missing directories), a \code{\textasciitilde} name is expanded, an absolute name is used
as given, and \code{-} keeps its meaning. A URL is not accepted for these files, also when
the input directory is a URL: the tool reports that it cannot read, or create, the file. In legacy mode, the input files are
opened as named, relative to the current directory, and the output directory, if one is
set, is a plain prefix of the names of the output files.

\subsection{Time arguments}
\label{sec:tools-times}

Options that take a time (\opt{-f}, \opt{-t}, \opt{-l} and others) accept the time strings
of the library function \code{strtim}:
\begin{itemize}
\item an elapsed time from the beginning of the record, \code{hh:mm:ss}, \code{mm:ss} or
  \code{ss}, with an optional fraction of a second (\code{1:30.5}, \code{10});
\item a sample number, \code{s} followed by the number (\code{s21600});
\item a time of day, in brackets: \code{[hh:mm:ss]}, \code{[hh:mm:ss dd/mm/yyyy]}, or
  \code{[hh:mm:ss d]} with a number of days after the start (it needs the base time of the
  record);
\item \code{e}, the end of the record; \code{c} followed by a counter value; \code{f}
  followed by a frame number.
\end{itemize}
Elapsed times and times of day need the sampling frequency of the record. When a record
does not give one, \wfdbx{} reports an error instead of assuming 250\,Hz
(\cref{chap:behavior}); sample numbers and \code{e} need no frequency. The annotation tools
\cmd{rdann}, \cmd{ann2rr}, \cmd{sortann}, \cmd{a2m} and \cmd{m2a} take \opt{-F}
\textit{FREQ} for that case, and \cmd{rr2ann} already had it.

\subsection{Help, options and arguments}
\label{sec:tools-help}

The options are those of the C tools: one letter after a hyphen, given separately
(\opt{-v -p}, not \opt{-vp}), their arguments as the next command-line arguments. As in
C, only the first letter after the hyphen is read in most tools, so that \opt{-pd} is
\opt{-p} with the time format \code{d}, and an argument that is not an option where an
option is expected is an error (\code{unrecognized argument}). Options that take a list
(\opt{-s} \textit{SIGNAL} \code{[}\textit{SIGNAL} \code{...]}, \opt{-p} \textit{TYPE} \code{...}) take
every following argument that does not begin with a hyphen.

\opt{-h} prints the tool's usage summary on the standard error. A tool that needs an
argument that was not given (often the record name) prints the same summary. The exit
status after \opt{-h} is that of the C tool, 0 or 1 depending on the tool (each entry says
which); as in C, \cmd{wfdbdesc}, \cmd{sampfreq} and \cmd{revise} take \opt{-h} alone as a
record name. The help texts of \cmd{ann2rr} and \cmd{wrsamp} pause after 23
lines, print \code{-{}-More-{}-} and read a line from the standard input, as the C tools do. The usage
line of each tool is given at the start of its entry below.

\subsection{Legacy mode}
\label{sec:tools-legacy}

\opt{-{}-legacy}, as the \emph{first} argument of a tool (\cmd{rdsamp -{}-legacy -r 100s},
\cmd{wfdbx rdann -{}-legacy -r 100s -a atr}), is removed before the options are read, and the
tool then behaves exactly as the original one: the options added by \wfdbx{} (such as
\opt{-F} of \cmd{rdann} or \opt{-R} of \cmd{wrsamp}) are unrecognized, the help texts and
messages are those of the original, the built-in values of the original (250\,Hz,
200\,adu/mV, ...) are used, files are found through the WFDB path (the environment variable
\code{WFDB}; the files that a tool reads itself are opened as named), and the exit status
is the original one. \opt{-{}-legacy} anywhere else is not
recognized: it is treated as any other argument. Legacy mode is meant for reproducing results of the original package;
\cref{chap:behavior} explains why it is not the default.

\subsection{Standard input and output}
\label{sec:tools-stdio}

Text output goes to the standard output, and messages, prompts and help texts to the
standard error. Some tools read the standard input (\cmd{wrsamp} without \opt{-i},
\cmd{edf2mit -i -}, \cmd{wrann}, \cmd{rr2ann}, \cmd{parsescp}, \cmd{readid}) and some write
binary data to the standard output (\cmd{wfdbcat}, \cmd{makeid}, \cmd{parsescp -a}, and
\cmd{wrsamp} without \opt{-o}).

\subsection{Exit status}
\label{sec:tools-status}

The exit status is that of the C tool: 0 on success, 1 for an error in the arguments
(including a required option that is missing), 2 when the record or a file cannot be read,
and, in some tools, other values that their entries give. In addition, \wfdbx{} follows the
convention of the standard Unix tools: when an output file or the standard output cannot be
written (a full disk, for example), a tool that would have exited with status 0 exits with
status 1, and reports a standard output that could not be written as
\code{PROGRAM: write error on standard output}, and a file that \cmd{snip}, \cmd{mit2edf}
or \cmd{parsescp} writes itself as \code{PROGRAM: write error on FILE}
(\cref{chap:behavior}). The original tools exited with status 0 in both cases; legacy
mode does too.

\section{Summary of the tools}
\label{sec:tools-summary}

\begin{xltabular}{\linewidth}{@{}P{2.3cm}L@{}}
\caption{The 33 tools of \wfdbx{}.}
\label{tab:tools}\\
\toprule
{\leavevmode\tabhead Tool} & {\leavevmode\tabhead What it does} \\
\midrule
\endfirsthead
\toprule
{\leavevmode\tabhead Tool} & {\leavevmode\tabhead What it does} \\
\midrule
\endhead
\bottomrule
\endfoot
\multicolumn{2}{@{}l}{\textit{Signal tools} (\cref{sec:tools-signal})} \\
\code{rdsamp} & print the samples of a record (text, CSV or XML; ADC or physical units) \\
\code{wrsamp} & write a record from a text file of samples \\
\code{wfdbdesc} & describe the signals of a record (the contents of its header) \\
\code{wfdbtime} & convert times to sample numbers, elapsed times and times of day \\
\code{time2sec} & convert a time to seconds \\
\code{sampfreq} & print the sampling frequency of a record \\
\code{signame} & print the names of the signals of a record \\
\code{signum} & print the numbers of the signals with given names \\
\code{sigamp} & measure signal amplitudes (peak-to-peak near beats, or RMS) \\
\code{snip} & copy a segment of a record to a new record \\
\code{xform} & convert a record: storage format, sampling frequency, gain, signals \\
\code{wfdbcat} & copy WFDB files to the standard output \\
\code{wfdbwhich} & print the path of a WFDB file \\
\midrule
\multicolumn{2}{@{}l}{\textit{Annotation tools} (\cref{sec:tools-annotation})} \\
\code{rdann} & print the annotations of an annotation file \\
\code{wrann} & write an annotation file from text in the format of \code{rdann} \\
\code{sortann} & sort an annotation file into canonical order \\
\code{ann2rr} & print the intervals between annotations (RR intervals) \\
\code{rr2ann} & write an annotation file from a list of RR intervals \\
\midrule
\multicolumn{2}{@{}l}{\textit{Format converters} (\cref{sec:tools-converters})} \\
\code{edf2mit} & convert an EDF or EDF+ file to a WFDB record \\
\code{mit2edf} & convert a WFDB record to an EDF or EDF+ file \\
\code{rdedfann} & print the annotations of an EDF+ file as \code{rdann} text \\
\code{wav2mit} & make a WAV file readable as a WFDB record \\
\code{mit2wav} & convert a WFDB record to a WAV file \\
\code{wfdb2mat} & convert a WFDB record to a MATLAB/Octave \file{.mat} file \\
\code{a2m} & convert an AHA-format annotation file to WFDB (MIT) format \\
\code{m2a} & convert a WFDB annotation file to AHA format \\
\code{ad2m} & convert an AHA-format signal file (tape or CD-ROM) to a WFDB record \\
\code{md2a} & convert a WFDB record to an AHA-format signal file \\
\code{ahaecg2mit} & convert AHA DVD files (binary or text) to WFDB records \\
\code{makeid} & write an AHA-format record ID block \\
\code{readid} & read an AHA-format record ID block \\
\code{revise} & convert an obsolete-format header to the current format \\
\code{parsescp} & read an SCP-ECG file, describe it, and convert it to a record \\
\end{xltabular}

% ==========================================================================================
\section{Signal tools}
\label{sec:tools-signal}

% ------------------------------------------------------------------------------------------
\apientry{rdsamp}{Print the samples of a record}
\begin{lstlisting}[style=signature,language={}]
usage: rdsamp -r RECORD [OPTIONS ...]
\end{lstlisting}
Reads the signals of a record and prints one line per sample interval: the sample number
(or, with \opt{-p}, the time) and the sample of each selected signal, in ADC units, or in
physical units with \opt{-p}. The output is plain text in columns (tab-separated), CSV
(\opt{-c}) or WFDB-XML (\opt{-X}), which \cmd{wrsamp} reads back. An invalid sample is
printed as \code{-} in physical units, and as the value $-32768$ in ADC units.
\begin{params}
\item[-r] The input record, \textit{RECORD} (required).
\item[-c] Use the CSV (comma-separated value) output format: the column headings of
  \opt{-v} are quoted (\code{'sample \#','MLII','V5'}), and the columns are separated by
  commas.
\item[-f] Begin at the time that follows (\cref{sec:tools-times}); by default, the beginning
  of the record.
\item[-h] Print the usage summary, and exit with status 0.
\item[-H] Read multifrequency records in high-resolution mode: one line per sample of the
  signals with the most samples per frame, the samples of the other signals repeated (by
  default, one line per frame, the samples of a frame averaged).
\item[-l] Stop after the time interval that follows (\code{hh:mm:ss}, or \code{s} followed
  by a number of samples), counted from the beginning time; a \opt{-t} time before the end
  of that interval still applies.
\item[-p] Print times and samples in physical units, with 3 decimal places (times in
  seconds by default). A letter appended to the option chooses the time column:
  \opt{-pd} the time of day and the date, when the record has a base time (elapsed time
  otherwise), \opt{-pe} the elapsed time as \code{hh:mm:ss.mmm}, \opt{-ph} hours,
  \opt{-pm} minutes, \opt{-ps} seconds (the default), \opt{-pS} the sample number.
\item[-P] As \opt{-p}, with 8 decimal places (columns 15 characters wide); \opt{-Pd},
  \opt{-Pe}, \opt{-Ph}, \opt{-Pm}, \opt{-Ps} and \opt{-PS} choose the time column.
\item[-s] Print only the signals that follow, given by number (0, 1, ...) or by name, in
  the order given: \opt{-s} \textit{SIGNAL} \code{[}\textit{SIGNAL} \code{...]}. A signal that
  the record does not have is an error.
\item[-S] Begin at the first valid sample of the signal that follows (a number or a name),
  at or after the \opt{-f} time.
\item[-t] Stop at the time that follows (by default, the end of the record).
\item[-v] Print column headings (the signal names, and with \opt{-p} the units).
\item[-X] Print WFDB-XML (an XML prolog, the sampling frequency, the number of signals,
  the headings and the samples); implies \opt{-c} and \opt{-v}.
\end{params}
The first four samples of \code{100s}, in ADC units, then four samples from 1 second in
physical units (\code{100s} gives no units, so they are shown as \code{nd}), and the CSV
format:
\begin{lstlisting}[style=text]
$ rdsamp -r 100s -t s4
              0     995    1011
              1     995    1011
              2     995    1011
              3     995    1011
$ rdsamp -r 100s -f 1 -t s364 -p -v
   Elapsed time    MLII      V5
      (seconds)    (nd)    (nd)
          1.000  -0.535  -0.205
          1.003  -0.505  -0.080
          1.006  -0.415   0.015
          1.008  -0.300   0.065
$ rdsamp -r 100s -t s3 -c -v
'sample #','MLII','V5'
0,995,1011
1,995,1011
2,995,1011
\end{lstlisting}
\returns{Exit status 0 when the samples have been printed (also when a read error or a
checksum error ends the output early, as in C).}
\raises{Exit status 1 for an error in the arguments (or a missing \opt{-r}); 2 if the record
cannot be opened, a signal of \opt{-s} or \opt{-S} cannot be found, or the sampling
frequency or a gain that the output needs is not specified.}
\textit{Differences from the C tool} (\cref{chap:behavior}): the sampling frequency is
required for time arguments other than sample numbers, for time columns (\opt{-p} without
\opt{-pS}, or with \opt{-v}) and for \opt{-X}; the gain of each selected signal is required
for physical units (\code{rdsamp: signal N of record R is uncalibrated}, status 2);
unspecified units are shown as \code{nd} (C: \code{mV}). The C tool assumed 250\,Hz and
200\,adu/mV. The quotes and backslashes of signal names are escaped in CSV headings, where
the C tool did not terminate.

% ------------------------------------------------------------------------------------------
\apientry{wrsamp}{Write a record from a text file of samples}
\begin{lstlisting}[style=signature,language={}]
usage: wrsamp [OPTIONS ...] [COLUMN ...]
\end{lstlisting}
Reads a text file with a sample of each signal on each line, and writes the samples to a
signal file and, with \opt{-o}, a header. Fields are separated by tabs, spaces, commas or
any combination of these (consecutive separators count as one). If the first line contains
letters, it gives the signal names; if the second line also contains letters, it gives the
units (with or without parentheses). WFDB-XML, as written by \cmd{rdsamp -X}, is read as
well. Non-numeric fields are written as invalid samples. Samples are rounded to the nearest
integer after scaling (\opt{-x}).
\begin{params}
\item[COLUMN] The fields to copy, by number (the leftmost field is column 0), in the order
  of the output signals; by default, all fields in order (see \opt{-z}).
\item[-c] Accepted for compatibility; it has no effect (as in the C tool).
\item[-d] Add dither (random noise of up to one unit) to the samples before rounding.
\item[-f] Start copying at the input line that follows (a number, $\ge 1$; the first line
  of samples is line 0).
\item[-F] The sampling frequency to write in the header (Hz). Required with \opt{-o}.
\item[-G] The gain (adu per physical unit) to write in the header; a quoted list of values
  (\opt{-G "100 50"}) gives one per signal, the last applying to the remaining signals. A
  gain of 0 declares the signals uncalibrated. Required with \opt{-o}.
\item[-h] Print the usage summary, and exit with status 1.
\item[-i] Read the input from the file that follows (by default, or with \code{-}, the
  standard input). The file is read from the input directory (\cref{sec:tools-files}).
\item[-o] Write the record that follows: the signal file \textit{RECORD}\file{.dat} and the
  header \textit{RECORD}\file{.hea}. Without \opt{-o}, the samples are written to the
  standard output (in format 16 unless \opt{-O} is given) and no header is written.
\item[-O] The storage format of the output (default 16); any format of
  \cref{tab:formats}.
\item[-r] The character that follows separates input lines (records), instead of the
  line feed.
\item[-R] The ADC resolution in bits; a quoted list gives one per signal, as for \opt{-G}.
  Required with \opt{-o}.
\item[-s] Any character of the string that follows separates input fields (by default:
  space, tab, carriage return or comma; only tab if the first line contains a tab).
\item[-t] Stop copying at the input line that follows (a number, $\ge 1$).
\item[-U] The units of the signals; a quoted list gives one per signal, as for \opt{-G}.
  Required with \opt{-o} for calibrated signals, unless the input gives the units.
\item[-x] Multiply the input values by the scale factor that follows; a quoted list gives
  one per signal.
\item[-z] Do not copy column 0 (usually a sample number), unless it is selected by a
  \textit{COLUMN} argument.
\end{params}
The obsolete option \opt{-l} (and its argument) is ignored, with a message, as in C. For
example, a record written from the output of \cmd{rdsamp} (column 0, the sample number, is
skipped by \opt{-z}):
\begin{lstlisting}[style=text]
$ rdsamp -r 100s -t s5 > s.txt
$ wrsamp -i s.txt -o rec -F 360 -G 200 -R 11 -U mV -z
$ cat rec.hea
rec 2 360 5
rec.dat 16 200/mV 11 0 995 4975 0 col 1
rec.dat 16 200/mV 11 0 1011 5055 0 col 2
$ wrsamp -i s.txt -o rec2
wrsamp: the sampling frequency must be specified (-F)
\end{lstlisting}
\returns{Exit status 0 when the record has been written.}
\raises{Exit status 1 for an error in the arguments, a missing \opt{-F}, \opt{-G}, \opt{-R}
or \opt{-U} (\code{wrsamp: the ... must be specified (-X)}), or an XML input without
samples; 2 if the input file cannot be read or the output cannot be opened; 3 if the input
contains no line separator.}
\textit{Differences from the C tool} (\cref{chap:behavior}): when a record is written
(\opt{-o}), the sampling frequency (\opt{-F}), the gains (\opt{-G}), the ADC resolutions
(the new option \opt{-R}) and, for calibrated signals, the units (the new option \opt{-U})
must be given; the C tool assumed 250\,Hz, 200\,adu/mV, a resolution based on the format,
and millivolts. A header written with \opt{-U mV} says \code{200/mV} where the C tool
wrote \code{200}. The \opt{-i} file is read from the input directory (C: from the current
directory).

% ------------------------------------------------------------------------------------------
\apientry{wfdbdesc}{Describe the signals of a record}
\begin{lstlisting}[style=signature,language={}]
usage: wfdbdesc RECORD [-readable]
\end{lstlisting}
Prints the contents of a record's header in readable form: the info strings (\code{Notes}),
the starting time, the length, the sampling frequency, and for each signal its file,
description, gain and units, initial value, storage format, ADC resolution and zero,
baseline and checksum. For a multi-segment record, it also gives the layout (fixed or
variable) and, for a fixed layout, the length of the first segment, whose signals it
describes.
\begin{params}
\item[RECORD] The record to describe.
\item[-readable] Describe only the signals whose files can be opened (by default, every
  signal named in the header, without opening the files). Any abbreviation of
  \opt{-readable} is accepted.
\end{params}
\begin{lstlisting}[style=text]
$ wfdbdesc 100s
Record 100s
Notes
=====
 69 M 1085 1629 x1
 Aldomet, Inderal
Produced by xform from record 100, beginning at 0:0
=====

Starting time: not specified
Length:     1:00.000 (21600 sample intervals)
Sampling frequency: 360 Hz
2 signals
Group 0, Signal 0:
 File: 100s.dat
 Description: MLII
 Gain: 200 adu/nd
 Initial value: 995
 Storage format: 212
 I/O: can be unbuffered
 ADC resolution: 11 bits
 ADC zero: 1024
 Baseline: 1024
 Checksum: 21537
Group 0, Signal 1:
 ...
\end{lstlisting}
\returns{Exit status 0.}
\raises{Exit status 1 without a record name; 2 if the header cannot be read, or the record
has no signals.}
\textit{Differences from the C tool} (\cref{chap:behavior}): an unspecified sampling
frequency or ADC resolution is shown as \code{not specified} (lengths are then given in
sample intervals only), a gain of 0 as \code{uncalibrated} (C:
\code{uncalibrated; assume 200}), and unspecified units as \code{adu/nd}; the C tool assumed
250\,Hz, 12 bits and millivolts. For a record whose header gives no length, the size of the
first signal file, from which the length is computed, is measured in the header's
directory, where the library reads the file (C: in the current directory). The notes are
the info strings that \code{getinfo} returns, after the record line found as the header
parser finds it (\cref{sec:record-line-divergence}): for a header whose first line is
empty, the comments before the record line are not printed (the C tool, and
\cmd{wfdbdesc -{}-legacy}, print them as notes).

% ------------------------------------------------------------------------------------------
\apientry{wfdbtime}{Convert times and sample numbers}
\begin{lstlisting}[style=signature,language={}]
usage: wfdbtime [ -r RECORD [ TIME ] ]
\end{lstlisting}
Converts each time argument, at the sampling frequency of the record, and prints three
columns: the sample number (\code{s} followed by the number), the elapsed time from the
beginning of the record, and the time of day (the elapsed time again if the record has no
base time).
\begin{params}
\item[-r] The record, \textit{RECORD}, whose sampling frequency and base time are used. A
  later \opt{-r} switches to another record for the times that follow it.
\item[TIME] A time to convert, in any form of \cref{sec:tools-times}; any number of times
  may follow the record.
\end{params}
\begin{lstlisting}[style=text]
$ wfdbtime -r 100s 1:00 s100
         s21600        1:00.000        1:00.000
           s100        0:00.278        0:00.278
\end{lstlisting}
\returns{Exit status 0.}
\raises{Exit status 1 without \opt{-r} (also for \opt{-h}, which prints the usage); 2 if
the record cannot be opened or its header gives no sampling frequency.}
\textit{Differences from the C tool} (\cref{chap:behavior}): the record must give its
sampling frequency (C: 250\,Hz assumed).

% ------------------------------------------------------------------------------------------
\apientry{time2sec}{Convert a time to seconds}
\begin{lstlisting}[style=signature,language={}]
usage: time2sec [-r RECORD] TIME
\end{lstlisting}
Prints a time interval in seconds. Without \opt{-r}, \textit{TIME} is converted at
1000 samples per second and printed with up to 12 significant digits; with \opt{-r}, it is
converted at the sampling frequency of the record and printed with at most three decimals
(milliseconds). A time of day in brackets (\code{[9:0:0]}) gives the elapsed time from the
beginning of the record.
\begin{params}
\item[-r] The record, \textit{RECORD}, whose sampling frequency (and base time) are used.
\item[TIME] The time to convert (\cref{sec:tools-times}).
\end{params}
\begin{lstlisting}[style=text]
$ time2sec 1:30.5
90.5
$ time2sec -r 100s s1000
2.778
\end{lstlisting}
\returns{Exit status 0.}
\raises{Exit status 1 for other arguments (also for \opt{-h}), after the usage summary.
Exit status 2 with \opt{-r} if the header of the record cannot be read, or gives no
sampling frequency (\code{time2sec: sampling frequency of record RECORD is not
specified}).}
\textit{Differences from the C tool}: with \opt{-r}, a record whose header gives no
sampling frequency (or cannot be read) is an error. The C tool converted the time at its
default frequency of 250\,Hz, a result unrelated to the record, because \wfdbx{} assumes no
frequency that a record does not give (\cref{sec:no-defaults}); until version 1.2.1,
\wfdbx{} divided by the unspecified frequency (0) and printed \code{-9223372036854775808}
with the exit status 0. With \opt{-{}-legacy}, the C behavior (250\,Hz).

% ------------------------------------------------------------------------------------------
\apientry{sampfreq}{Print the sampling frequency of a record}
\begin{lstlisting}[style=signature,language={}]
usage: sampfreq [OPTION] RECORD
\end{lstlisting}
Prints the sampling frequency (frame rate) of a record, in Hz. With \opt{-a}, prints the
name and sampling frequency of each signal, tab-separated; with \opt{-H}, the highest
sampling frequency of a multifrequency record.
\begin{params}
\item[RECORD] The record.
\item[-a] List all signals and their sampling frequencies.
\item[-h] Listed in the usage summary, which is printed (status 1) when no argument is
  given; as in C, \cmd{sampfreq -h} alone takes \code{-h} as a record name.
\item[-H] Print the maximum sampling frequency of a multifrequency record.
\end{params}
\begin{lstlisting}[style=text]
$ sampfreq 100s
360
$ sampfreq -a mfreq
slow    100
medium  200
fast    400
$ sampfreq -H mfreq
400
\end{lstlisting}
\returns{Exit status 0; a header that cannot be read prints the library's message and
\code{-1}.}
\raises{Exit status 1 without arguments.}
\textit{Differences from the C tool} (\cref{chap:behavior}): for a header that does not
give a sampling frequency, the library reports 0, which is printed (C: 250).

% ------------------------------------------------------------------------------------------
\apientry{signame}{Print the names of the signals of a record}
\begin{lstlisting}[style=signature,language={}]
usage: signame -r RECORD [OPTIONS ...]
\end{lstlisting}
Prints the name (description) of each signal of the record, one per line, or of the signals
whose numbers follow \opt{-s}. A number that is not that of a signal prints
\code{[INVALID]}.
\begin{params}
\item[-r] The record, \textit{RECORD} (required).
\item[-h] Print the usage summary, and exit with status 0.
\item[-s] Print the names of the signals whose numbers follow:
  \opt{-s} \textit{SIGNAL} \code{[}\textit{SIGNAL} \code{...]}.
\end{params}
\begin{lstlisting}[style=text]
$ signame -r 100s -s 1 0 2
V5
MLII
[INVALID]
\end{lstlisting}
\returns{Exit status 0.}
\raises{Exit status 1 for an error in the arguments; 2 if the record cannot be opened or
has no signals.}
\textit{Differences from the C tool}: none, apart from \opt{-{}-legacy} (the C tool also
read past the end of its array for the number equal to the number of signals; \wfdbx{}
prints \code{[INVALID]}).

% ------------------------------------------------------------------------------------------
\apientry{signum}{Print the numbers of the signals with given names}
\begin{lstlisting}[style=signature,language={}]
usage: signum -r RECORD [OPTIONS ...]
\end{lstlisting}
For each name that follows \opt{-s}, prints the numbers of the signals of the record with
that exact name (tab-separated, when several have it), or \code{X} if none has it, one line
per name.
\begin{params}
\item[-r] The record, \textit{RECORD} (required).
\item[-h] Print the usage summary, and exit with status 0.
\item[-s] The signal names to look for (required):
  \opt{-s} \textit{SIGNAL} \code{[}\textit{SIGNAL} \code{...]}.
\end{params}
\begin{lstlisting}[style=text]
$ signum -r 100s -s V5 MLII ABP
1
0
X
\end{lstlisting}
\returns{Exit status 0.}
\raises{Exit status 1 for an error in the arguments, or without \opt{-r} or \opt{-s};
2 if the record cannot be opened.}
\textit{Differences from the C tool}: none, apart from \opt{-{}-legacy}.

% ------------------------------------------------------------------------------------------
\apientry{sigamp}{Measure signal amplitudes}
\begin{lstlisting}[style=signature,language={}]
usage: sigamp -r RECORD [OPTIONS ...]
\end{lstlisting}
Measures the amplitude of each signal, up to \textit{NMAX} times, and prints the trimmed
mean of the measurements (the largest and smallest 5\% excluded), one tab-separated column
per signal. With an annotator (\opt{-a}), each measurement is the peak-to-peak amplitude in
a window around a beat annotation (an annotation that \code{map1} maps to \code{NORMAL});
without one, it is the RMS amplitude over consecutive windows of 1 second (\opt{-w}). The
amplitudes are in ADC units, or in physical units with \opt{-p}.
\begin{params}
\item[-r] The record, \textit{RECORD} (required).
\item[-a] Measure peak-to-peak amplitudes near the beat annotations of the annotator that
  follows. \opt{-a} and \opt{-w} cannot be used together (\opt{-w} is then ignored, with a
  message).
\item[-d] The window of the peak-to-peak measurements, \textit{DT1} \textit{DT2}: the time
  before and the time after each annotation (default 0.05 seconds each); used only with
  \opt{-a}.
\item[-f] Begin at the time that follows.
\item[-h] Print the usage summary, and exit with status 0.
\item[-H] Read multifrequency records in high-resolution mode.
\item[-n] Make up to the number of measurements that follows per signal (default 300).
\item[-p] Print the amplitudes in physical units. A letter appended chooses the time
  format of the individual measurements (\opt{-q}, \opt{-v}): \opt{-pd} the time of day,
  \opt{-pe} \code{hh:mm:ss}, \opt{-ph} hours, \opt{-pm} minutes, \opt{-ps} seconds (the
  default), \opt{-pS} sample numbers.
\item[-q] Quick mode: print the individual measurements only, each with its time.
\item[-t] Stop at the time that follows.
\item[-v] Verbose mode: print the individual measurements, then the trimmed means
  (\code{Trimmed mean}).
\item[-w] The window of the RMS measurements, \textit{DTW} (default 1 second).
\end{params}
The environment variable \code{WFDBGVMODE} sets the resolution mode when \opt{-H} is not
given, as in C.
\begin{lstlisting}[style=text]
$ sigamp -r 100s
        32.97   23.7155
$ sigamp -r 100s -a atr -n 4 -v
0.213889        265     176
1.027778        295     193
1.838889        306     216
2.627778        301     216
Trimmed mean    291.75  200.25
\end{lstlisting}
\returns{Exit status 0.}
\raises{Exit status 1 for an error in the arguments; 2 if the record or the annotator cannot
be opened, the record gives no sampling frequency, or a signal is uncalibrated with
\opt{-p}.}
\textit{Differences from the C tool} (\cref{chap:behavior}): the record must give its
sampling frequency (the windows are in seconds), and, with \opt{-p}, the gain of each signal
(C: 250\,Hz and 200\,adu/mV assumed).

% ------------------------------------------------------------------------------------------
\apientry{snip}{Copy a segment of a record to a new record}
\begin{lstlisting}[style=signature,language={}]
usage: snip -i IREC -n NREC [OPTIONS ...]
\end{lstlisting}
Copies the samples of a record between two times to a new record, with the same signals,
storage format and signal specifications, and optionally its annotations, with times
relative to the beginning of the excerpt. The header of the new record keeps the info
strings of the input and adds \code{Produced by snip from record IREC, beginning at TIME}.
To change the sampling frequency, gains or signals, \cmd{xform} is the tool to use.
\begin{params}
\item[-i] The input record, \textit{IREC} (required; \opt{-r} is accepted as a synonym).
\item[-n] The new record, \textit{NREC} (required). It must differ from \textit{IREC}: the
  header of the new record is written over that of a record of the same name.
\item[-a] Copy the annotations of the annotators that follow:
  \opt{-a} \textit{ANNOTATOR} \code{[}\textit{ANNOTATOR} \code{...]}.
\item[-f] Begin at the time that follows (default 0).
\item[-h] Print the usage summary, and exit with status 0.
\item[-l] Copy only the duration that follows (\code{hh:mm:ss} or \code{s} followed by a
  number of samples); overrides \opt{-t}.
\item[-m] Keep the segments of a multi-segment input record, if possible. The master
  header of the new record is created in the output directory, as the other files are.
\item[-O] Write the samples in the storage format that follows (by default, the format of
  the input signals, or, if they differ, the best fit to the largest ADC resolution).
\item[-s] Do not write info strings in the new header (neither those of the input nor the
  \code{Produced by snip} line).
\item[-t] Stop at the time that follows (default: the end of the record).
\end{params}
\begin{lstlisting}[style=text]
$ snip -i 100s -n ex -f 10 -t 20 -a atr
$ cat ex.hea
ex 2 360 3600
ex.dat 212 200 11 1024 946 -16262 0 MLII
ex.dat 212 200 11 1024 969 18142 0 V5
# 69 M 1085 1629 x1
# Aldomet, Inderal
#Produced by xform from record 100, beginning at 0:0
#Produced by snip from record 100s, beginning at 10
\end{lstlisting}
\returns{Exit status 0.}
\raises{Exit status 1 for an error in the arguments or an improper interval; 2 if a header
cannot be read or written, or the input record gives no sampling frequency.}
\textit{Differences from the C tool} (\cref{chap:behavior}): the input record must give its
sampling frequency, for the new header (C: 250\,Hz assumed); an output signal file is never
one of the input files; when the last frame to copy cannot be read and no frame was read
before it, 0 is written (C: uninitialized values); the master header of \opt{-m} is created
in the output directory (C: in the current directory).

% ------------------------------------------------------------------------------------------
\apientry{xform}{Convert a record (format, frequency, gain, signals)}
\begin{lstlisting}[style=signature,language={}]
usage: xform -i IREC [OPTIONS ...]
\end{lstlisting}
Copies the signals of a record to a new one, converting them: to another storage format,
sampling frequency (by interpolation), gain, ADC resolution or zero level, and to a subset
or another order of the signals; annotations can be copied too, with their times converted.
The output is described either by an existing header (\opt{-o}) or by answers to prompts
(printed on the standard error and read from the terminal, the standard input, or a script
given by \opt{-S}): the output record name (unless \opt{-n} gives it), the number of
signals, the sampling frequency, the directory of the signal files, the storage format,
whether all signals go in one file, and for each signal its gain, ADC resolution and ADC
zero (and with \opt{-N} its description and units). Samples out of the range of the output
format wrap around unless \opt{-c} is given. While it converts, it prints
\code{Checking input signals ...} and progress dots.
\begin{params}
\item[-i] The input record, \textit{IREC} (required).
\item[-a] Copy the annotations of the annotators that follow:
  \opt{-a} \textit{ANNOTATOR} \code{[}\textit{ANNOTATOR} \code{...]}.
\item[-c] Clip the output samples to the range of the format (by default, they wrap around).
\item[-d] Add dither to the input if the sampling frequency or the gain changes.
\item[-f] Begin at the time that follows.
\item[-h] Print the usage summary, and exit with status 0.
\item[-H] Open the input record in high-resolution mode.
\item[-M] Process a multifrequency record without changing the frequencies of its signals.
\item[-n] Create a header for the record \textit{NREC} that follows, with the signal
  descriptions of \textit{IREC}.
\item[-N] As \opt{-n}, but with the signal descriptions of \textit{OREC} (or asked for).
\item[-o] Write the output signal files as the existing header of the record \textit{OREC}
  that follows describes them (no prompts).
\item[-s] Write only the input signals that follow, by number or name, in that order:
  \opt{-s} \textit{SIGNAL} \code{[}\textit{SIGNAL} \code{...]}. A signal that the record does not
  have is skipped, with a warning.
\item[-S] Take the answers to the prompts from the text file \textit{SCRIPT} that follows,
  one answer per line (an empty line takes the default answer). The file is read from the
  input directory (\cref{sec:tools-files}).
\item[-t] Stop at the time that follows.
\item[-u] Adjust annotation times, if needed, to make them unique.
\end{params}
For example, a copy of the first 2 seconds of \code{100s} in format 16, with the answers in
the file \file{ans.txt} (empty lines, except \code{16} on the fourth: the default number
of signals, frequency and directory, then format 16, then the default answers):
\begin{lstlisting}[style=text]
$ xform -i 100s -n x16 -S ans.txt -t 2 2>/dev/null
$ cat x16.hea
x16 2 360 720
x16.dat 16 200 11 1024 995 -26770 0 MLII
x16.dat 16 200 11 1024 1011 -12260 0 V5
# 69 M 1085 1629 x1
# Aldomet, Inderal
#Produced by xform from record 100, beginning at 0:0
#Produced by xform from record 100s, beginning at 0:0
\end{lstlisting}
\returns{Exit status 0.}
\raises{Exit status 1 for an error in the arguments, identical \opt{-n} and \opt{-o} names,
\opt{-M} with a change of frequency, an output frequency of 0, a missing ADC resolution, a
\opt{-t} time before the \opt{-f} time, or the end of a script before the last answer; 2 if
a record cannot be opened or written, a sampling frequency is not specified, or a requested
rescaling is impossible.}
\textit{Differences from the C tool} (\cref{chap:behavior}): the input record (and the
record of \opt{-o}) must give a sampling frequency, and an output frequency of 0 is refused
(C: 250\,Hz); rescaling an uncalibrated signal to a given gain, and rescaling by ADC
resolution when a resolution is unspecified, are refused; an ADC resolution is offered as
the default only if it comes from the input (or the previous output signal) and suits the
output format, and must be given otherwise; unspecified units are shown as \code{nd} in the
prompts; an output signal file that would be one of the input files is refused
(\code{osigfopen: output file NAME is also an input file}). With \opt{-M} and an output
frame larger than the input frame, the samples beyond the input frame are 0 (C: a crash).
The script of \opt{-S} is read from the input directory (C: from the current directory).

% ------------------------------------------------------------------------------------------
\apientry{wfdbcat}{Copy WFDB files to the standard output}
\begin{lstlisting}[style=signature,language={}]
usage: wfdbcat FILE [FILE ...]
\end{lstlisting}
Copies each file, found as the library finds files (\cref{sec:tools-files}: relative to the
input directory, or as a full path or a URL), to the standard output, byte for byte. A file
that cannot be found is reported (\code{wfdbcat: `FILE' not found}) and skipped.
\begin{params}
\item[FILE] A file to copy, named with its suffix (\code{100s.hea}).
\end{params}
\begin{lstlisting}[style=text]
$ wfdbcat 100s.hea
100s 2 360 21600
100s.dat 212 200 11 1024 995 21537 0 MLII
100s.dat 212 200 11 1024 1011 -3962 0 V5
# 69 M 1085 1629 x1
# Aldomet, Inderal
#Produced by xform from record 100, beginning at 0:0
\end{lstlisting}
\returns{Exit status 0, also when a file was not found.}
\raises{Exit status 1 without arguments, or with \opt{-h} alone.}
\textit{Differences from the C tool} (\cref{chap:behavior}): each file is read from where it
is named, not found by a search of the WFDB path.

% ------------------------------------------------------------------------------------------
\apientry{wfdbwhich}{Print the path of a WFDB file}
\begin{lstlisting}[style=signature,language={}]
usage: wfdbwhich [-r RECORD ] FILENAME
 -or-  wfdbwhich FILE-TYPE RECORD
\end{lstlisting}
Prints the path by which the library opens a file: a file named by \textit{FILENAME}, or
the file of type \textit{FILE-TYPE} (a suffix such as \code{hea}, \code{dat} or \code{atr})
of a record. A file that cannot be found is reported with the input directory
(\code{(WFDB path = ...)}).
\begin{params}
\item[-r] Open the record \textit{RECORD} first.
\item[FILENAME] The name of the file to find.
\item[FILE-TYPE] The type (suffix) of the file of the record to find.
\item[RECORD] The record.
\end{params}
\begin{lstlisting}[style=text]
$ wfdbwhich hea 100s
100s.hea
$ wfdbwhich 100x.dat
`100x.dat' not found
(WFDB path = .)
\end{lstlisting}
\returns{Exit status 0 when the file is found.}
\raises{Exit status 1 for other arguments (also \opt{-h}); 2 if the file is not found.}
\textit{Differences from the C tool} (\cref{chap:behavior}): the file is looked for only
where it is named (the input directory, a full path or a URL), not by a search of the WFDB
path, and the path printed in the message is the input directory.

% ==========================================================================================
\section{Annotation tools}
\label{sec:tools-annotation}

% ------------------------------------------------------------------------------------------
\apientry{rdann}{Print the annotations of an annotation file}
\begin{lstlisting}[style=signature,language={}]
usage: rdann -r RECORD -a ANNOTATOR [OPTIONS ...]
\end{lstlisting}
Prints one line per annotation: the time (elapsed time, or time of day in brackets when the
record has a base time), the sample number, the mnemonic of the annotation type, the
subtype, channel and number fields, and the auxiliary information (after a tab), if any.
Options select the annotations to print and the time format. This text is the input format
of \cmd{wrann}. The old syntax of the C tool, \cmd{rdann} \textit{ANNOTATOR RECORD}
\code{[}\textit{FROM} \code{[}\textit{TO} \code{[}\textit{TYPE} \code{...]]]}, is accepted too.
\begin{params}
\item[-r] The record, \textit{RECORD} (required).
\item[-a] The annotator, \textit{ANNOTATOR} (required).
\item[-c] Print only the annotations with the channel number that follows (from $-128$
  to 127).
\item[-e] Print the times as elapsed times, even for a record with a base time. Only one of
  \opt{-e} and \opt{-x} can be used (the last one given applies).
\item[-f] Start at the time that follows; also \code{a} followed by a number (from that
  annotation number), or \code{b} or \code{\#} followed by a number (from that beat number).
\item[-F] The sampling frequency, used if the record has no header or its header gives
  none.
\item[-h] Print the usage summary, and exit with status 0.
\item[-n] Print only the annotations with the \code{num} value that follows.
\item[-p] Print only the annotations of the types that follow, given by their mnemonics:
  \opt{-p} \textit{TYPE} \code{[}\textit{TYPE} \code{...]} (\opt{-p N V}).
\item[-s] Print only the annotations with the subtype that follows.
\item[-t] Stop at the time that follows; also \code{a}, \code{b} or \code{\#} followed by a
  number, as for \opt{-f} (\code{\#}: that number of beats).
\item[-v] Print column headings.
\item[-x] Print the times in seconds, minutes and hours (with enough decimal places to
  resolve the time resolution of the file).
\end{params}
\begin{lstlisting}[style=text]
$ rdann -r 100s -a atr -t 3 -v
      Time   Sample #  Type  Sub Chan  Num      Aux
    0:00.050       18     +    0    0    0      (N
    0:00.214       77     N    0    0    0
    0:01.028      370     N    0    0    0
    0:01.839      662     N    0    0    0
    0:02.628      946     N    0    0    0
$ rdann -r 100s -a atr -t 2 -x
    0.050   0.00083 0.0000139     +    0    0    0      (N
    0.214   0.00356 0.0000594     N    0    0    0
    1.028   0.01713 0.0002855     N    0    0    0
    1.839   0.03065 0.0005108     N    0    0    0
\end{lstlisting}
\returns{Exit status 0.}
\raises{Exit status 1 for an error in the arguments, or if the sampling frequency is not
specified (\code{rdann: sampling frequency of record R is not specified (use -F)}); 2 if the
annotation file cannot be opened or a starting annotation or beat number is beyond its end.}
\textit{Differences from the C tool} (\cref{chap:behavior}): the new option \opt{-F} gives
the sampling frequency of a record without a header (or whose header gives none), and is
then required; the C tool assumed 250\,Hz.

% ------------------------------------------------------------------------------------------
\apientry{wrann}{Write an annotation file from text}
\begin{lstlisting}[style=signature,language={}]
usage: wrann -r RECORD -a ANNOTATOR <TEXT-FILE
\end{lstlisting}
Reads annotations from the standard input, in the format printed by \cmd{rdann}, and writes
the annotation file \textit{RECORD}\file{.}\textit{ANNOTATOR}. The time column is skipped and
the sample number is used, except on lines that begin with a time of day in brackets, whose
time is converted. A mnemonic that is not an annotation code is written as a \code{NOTE}
(comment) annotation whose auxiliary information begins with that mnemonic. Text after a
tab is the auxiliary information. The time resolution recorded in the file is the sampling
frequency of the record, if its header gives one. Annotations that are not in canonical
order are sorted when the file is closed (\cref{chap:behavior}).
\begin{params}
\item[-r] The record, \textit{RECORD} (required).
\item[-a] The output annotator, \textit{ANNOTATOR} (required).
\end{params}
\opt{-h}, not listed in the usage, prints it and exits with status 0.
\begin{lstlisting}[style=text]
$ rdann -r 100s -a atr | wrann -r 100s -a copy
\end{lstlisting}
\returns{Exit status 0.}
\raises{Exit status 1 for an error in the arguments; 2 if the output file cannot be opened.}
\textit{Differences from the C tool} (\cref{chap:behavior}): an out-of-order file is sorted
by the library itself, without losing annotations, instead of by running \cmd{sortann}.

% ------------------------------------------------------------------------------------------
\apientry{sortann}{Sort an annotation file}
\begin{lstlisting}[style=signature,language={}]
usage: sortann -r RECORD -a ANNOTATOR [OPTIONS ...]
\end{lstlisting}
Rearranges the annotations of a file in canonical order: by time, then \code{num}, then
\code{chan}. A later annotation with the same time, \code{num} and \code{chan} replaces an
earlier one, and \code{NOTQRS} annotations are dropped. Without \opt{-o}, the sorted file is
written under the name of the input (in the output directory, so that it replaces the
input when the input and output directories are the same), unless the input is already
sorted: that is reported (\code{sortann: input is already ordered -{}- no output written})
and nothing is written.
\begin{params}
\item[-r] The record, \textit{RECORD} (required). As in C, \opt{-r}\textit{DIR} (the
  directory attached to the option) sets the input directory before the record name.
\item[-a] The input annotator, \textit{ANNOTATOR} (required).
\item[-f] Copy only the annotations from the time that follows (requires \opt{-o}).
\item[-F] The sampling frequency, used if the record has no header or its header gives
  none; then required for \opt{-f} and \opt{-t}.
\item[-h] Print the usage summary, and exit with status 0.
\item[-o] Write the output as the annotator \textit{OUTANN} that follows (it must differ
  from the input annotator); required with \opt{-f} or \opt{-t}.
\item[-t] Copy only the annotations before the time that follows (requires \opt{-o}).
\end{params}
\begin{lstlisting}[style=text]
$ sortann -r 100s -a atr
sortann: input is already ordered -- no output written
\end{lstlisting}
\returns{Exit status 0, also for an input that is already sorted.}
\raises{Exit status 1 for an error in the arguments, a missing \opt{-o}, identical input
and output annotators, or a sampling frequency that is needed and not specified; 2 if an
annotation file cannot be opened.}
\textit{Differences from the C tool} (\cref{chap:behavior}): the new option \opt{-F}; an
out-of-order annotation at time 0 (with number and channel 0) is no longer lost, nor are
those inserted after it (the C tool dropped them; legacy mode does too).

% ------------------------------------------------------------------------------------------
\apientry{ann2rr}{Print the intervals between annotations}
\begin{lstlisting}[style=signature,language={}]
usage: ann2rr -r RECORD -a ANNOTATOR [OPTIONS ...]
\end{lstlisting}
Prints the intervals between consecutive beat (QRS) annotations, one per line, in sample
intervals by default, optionally with the times and types of the annotations that begin and
end each interval. The first interval begins at the starting time.
\begin{params}
\item[-r] The record, \textit{RECORD} (required).
\item[-a] The annotator, \textit{ANNOTATOR} (required).
\item[-A] Print all intervals between annotations, not only RR intervals (overrides \opt{-c}
  and \opt{-p}).
\item[-c] Print only intervals between consecutive valid annotations.
\item[-f] Start at the time that follows.
\item[-F] The sampling frequency, used if the record has no header or its header gives
  none; then required for times and for intervals in units of time.
\item[-h] Print the usage summary, and exit with status 0.
\item[-i] Print the intervals in the format that follows: \code{h} (hours), \code{m}
  (minutes), \code{s} (seconds), each optionally followed by a number of decimal places from
  0 to 15 (default 3; \opt{-i s8}), or \code{t} (\code{hh:mm:ss}).
\item[-p] Print only intervals that end with annotations of the types that follow:
  \opt{-p} \textit{TYPE} \code{[}\textit{TYPE} \code{...]}, given by mnemonics such as
  \code{N} or \code{V}.
\item[-P] Print only intervals that begin with annotations of the types that follow.
\item[-t] Stop at the time that follows.
\item[-v] Print the time of the end of each interval, in the format that follows, if any:
  \code{h}, \code{m} or \code{s} (with decimal places, as for \opt{-i}), \code{t}, or
  \code{T} (date and time, if available); without a format, in sample intervals.
\item[-V] Print the time of the beginning of each interval, as for \opt{-v}.
\item[-w] Print the annotations that end the intervals.
\item[-W] Print the annotations that begin the intervals.
\end{params}
\begin{lstlisting}[style=text]
$ ann2rr -r 100s -a atr -t 4 -i s -V s -w
0.000   0.214   N
0.214   0.814   N
1.028   0.811   N
1.839   0.789   N
2.628   0.792   N
\end{lstlisting}
\returns{Exit status 0.}
\raises{Exit status 1 for an error in the arguments, or a sampling frequency that is needed
and not specified; 2 if the annotation file cannot be opened.}
\textit{Differences from the C tool} (\cref{chap:behavior}): the new option \opt{-F}; the C
tool assumed 250\,Hz.

% ------------------------------------------------------------------------------------------
\apientry{rr2ann}{Write an annotation file from RR intervals}
\begin{lstlisting}[style=signature,language={}]
usage: rr2ann -r RECORD -a ANNOTATOR [OPTIONS ...] <TEXT-FILE
\end{lstlisting}
Reads a column of RR intervals (or of times of occurrence, with \opt{-T}) in sample
intervals from the standard input, and writes an annotation file of \code{N} (normal beat)
annotations at the corresponding times, or of the types of a second column (\opt{-w}). The
first token of a line is the interval; empty lines, non-numeric tokens, and negative or zero
intervals are skipped. If the record has no header, one is created.
\begin{params}
\item[-r] The record, \textit{RECORD} (required).
\item[-a] The output annotator, \textit{ANNOTATOR} (required).
\item[-F] The sampling frequency, used (and required) if the header does not exist or does
  not give one; otherwise ignored, with a warning.
\item[-h] Print the usage summary, and exit with status 0.
\item[-T] The input contains times of occurrence, not RR intervals.
\item[-x] Multiply the input values by the factor that follows (default 1) to obtain sample
  intervals.
\item[-w] Take the annotation type of each line from its second column (a mnemonic).
\end{params}
\begin{lstlisting}[style=text]
$ ann2rr -r 100s -a atr | rr2ann -r 100s -a rr
$ rdann -r 100s -a rr -t 2
    0:00.214       77     N    0    0    0
    0:01.028      370     N    0    0    0
    0:01.839      662     N    0    0    0
\end{lstlisting}
\returns{Exit status 0.}
\raises{Exit status 1 for an error in the arguments, or a sampling frequency that is needed
and not given; 2 if the header cannot be created; 3 if the annotation file cannot be
opened.}
\textit{Differences from the C tool} (\cref{chap:behavior}): \opt{-F} is required when the
record has no header or its header gives no frequency (C: 250\,Hz assumed); a header is
created only when there is none.

% ==========================================================================================
\section{Format converters}
\label{sec:tools-converters}

% ------------------------------------------------------------------------------------------
\apientry{edf2mit}{Convert an EDF or EDF+ file to a WFDB record}
\begin{lstlisting}[style=signature,language={}]
usage: edf2mit [OPTIONS ...]
\end{lstlisting}
Reads a European Data Format file (EDF or EDF+) and writes a WFDB record with the same
data: a header and a signal file in format 16, with the sampling frequencies, gains,
baselines, units and start time of the EDF header. The signals \code{EDF Annotations} of an
EDF+ file are copied as signals (\cmd{rdedfann} reads them). Unknown options are ignored.
\begin{params}
\item[-i] The EDF file to read, \textit{EDFILE} (required); \code{-} reads the standard
  input. It is read from the input directory (\cref{sec:tools-files}).
\item[-b] The input is in big-endian byte order (by default, little-endian).
\item[-h] Print the usage summary, and exit with status 0.
\item[-r] The record to create, \textit{RECORD} (by default, the input file name up to its
  first dot, with every character other than a letter or digit replaced by \code{\_}).
\item[-s] Copy only the signals that follow, by number (from 0):
  \opt{-s} \textit{SIGNAL} \code{[}\textit{SIGNAL} \code{...]}.
\item[-v] Print the fields of the EDF header (debugging output) on the standard output.
\end{params}
\begin{lstlisting}[style=text]
$ edf2mit -i edfp.edf -r edfp2
$ cat edfp2.hea
edfp2 2 10 20 13:24:35 02/03/2002
edfp2.dat 16x10 4.095/uV 11 0 0 26823 0 EEG Fpz-Cz
edfp2.dat 16x3 32767.5 15 0 12331 14629 0 EDF Annotations
\end{lstlisting}
\begin{phibox}
The patient and recording identification fields of an EDF header can identify a person;
\opt{-v} prints them.
\end{phibox}
\returns{Exit status 0.}
\raises{Exit status 1 for a missing or unreadable input file or an invalid record name;
2 if the input is not EDF, or is shorter than the 256-byte fixed part of an EDF header
(\code{edf2mit: incomplete EDF header in FILE}; no file is created).}
\textit{Differences from the C tool} (\cref{chap:behavior}): the \opt{-i} file is read from
the input directory (C: from the current directory); otherwise none, apart from
\opt{-{}-legacy}, except for input on which the C tool fails: a file shorter than the fixed
part of the header (C: fields from uninitialized memory), and signals with 0 samples per
data record (C: a crash).

% ------------------------------------------------------------------------------------------
\apientry{mit2edf}{Convert a WFDB record to an EDF file}
\begin{lstlisting}[style=signature,language={}]
usage: mit2edf -r RECORD [OPTIONS ...]
\end{lstlisting}
Reads a record and writes an EDF file with the same data: 16-bit samples in data records
(blocks) of 10 seconds, or shorter if a block would exceed 61440 bytes. The record name is
written as the patient identification, and the base date and time, if any, as the start
date and time (\code{Startdate not recorded} otherwise). An empty line is printed on the
standard output when the file is complete. Unknown options are ignored.
\begin{params}
\item[-r] The input record, \textit{RECORD} (required).
\item[-h] Print the usage summary, and exit with status 0.
\item[-o] The EDF file to write, \textit{EDFILE} (default: \textit{RECORD}\file{.edf},
  without the directory part of the record name), created in the output directory
  (\cref{sec:tools-files}).
\item[-p] Write an EDF+ file (its name must end in \file{.edf} or \file{.EDF}); warnings say
  what EDF+ requires that the record does not give.
\item[-v] Verbose mode: print the sizes of the header, the blocks and the file, and a dot
  per minute of data.
\end{params}
\begin{lstlisting}[style=text]
$ mit2edf -r 100s -o 100s.edf -v
Converting record 100s to 100s.edf (EDF mode)
 Header block size: 768 bytes
 Data block size: 10 seconds (3600 frames or 14400 bytes)
 Recording length: 1:00 (7 data blocks, 25200 frames, 100800 bytes)
 Total length of file to be written: 101568 bytes
.
\end{lstlisting}
\returns{Exit status 0.}
\raises{Exit status 1 for an error in the arguments or an invalid EDF+ file name; 2 if the
record cannot be opened; 3 if its signals cannot be read, or the sampling frequency, a gain
or an ADC resolution is not specified; 4 if the output file cannot be created; 5 if a frame
does not fit in an EDF block.}
\textit{Differences from the C tool} (\cref{chap:behavior}): the record must give the
sampling frequency and the gain and ADC resolution of every signal (an EDF file stores
physical and digital ranges); the C tool assumed 250\,Hz and 200\,adu/mV and guessed the
resolution from the storage format. Unspecified units are written as \code{nd} (C:
\code{mV}). The EDF file is created in the output directory (C: in the current directory).

% ------------------------------------------------------------------------------------------
\apientry{rdedfann}{Print the annotations of an EDF+ file}
\begin{lstlisting}[style=signature,language={}]
usage: rdedfann -r RECORD [OPTIONS ...]
\end{lstlisting}
Reads the \code{EDF Annotations} signal of an EDF+ file, read as a record
(\cref{chap:records}), and prints its annotations as text in the format of \cmd{rdann},
which \cmd{wrann} can convert to an annotation file. The annotation text is printed in place
of the mnemonic (by default) or as the auxiliary information (\opt{-x}), followed by the
duration, if any.
\begin{params}
\item[-r] The EDF+ file, read as a record, \textit{RECORD} (required): \opt{-r edfp.edf}.
\item[-F] Compute the sample numbers at the frequency \textit{FREQ} that follows, in Hz
  (a value that is not positive gives 1).
\item[-h] Print the usage summary, and exit with status 0.
\item[-v] Print column headings.
\item[-x] Print the annotation text in the auxiliary information, with the type
  \code{\textquotedbl} (\code{NOTE}).
\end{params}
\begin{lstlisting}[style=text]
$ rdedfann -r edfp.edf -v
      Time   Sample #  Type  Sub Chan  Num      Aux
    0:00.500       50 Sleep_stage_W    0    0    0              duration: 1.5
    0:01.250      125 Arousal    0    0    0            duration: 1.5
    0:01.250      125 Spindle_onset    0    0    0              duration: 1.5
\end{lstlisting}
\returns{Exit status 0.}
\raises{Exit status 1 for an error in the arguments; 2 if the record cannot be opened; 3 if
it has no EDF annotations.}
\textit{Differences from the C tool}: none, apart from \opt{-{}-legacy}.

% ------------------------------------------------------------------------------------------
\apientry{wav2mit}{Make a WAV file readable as a record}
\begin{lstlisting}[style=signature,language={}]
usage: wav2mit -i FILE.wav -r RECORD [ OPTIONS ... ]
\end{lstlisting}
Writes a header for a WAV file (PCM, 8 or 16 bits per sample), so that the WAV file itself
is read as the signal file of a record: format 16 (or 80 for 8-bit files) after a 44-byte
prolog, with the sampling frequency, number of signals and resolution of the WAV file.
\begin{params}
\item[-i] The WAV file, \textit{FILE.wav} (required; the name must end in \file{.wav}),
  read from the input directory (\cref{sec:tools-files}).
\item[-G] The gain in adu per physical unit (required; 0 declares the signals
  uncalibrated).
\item[-h] Print the usage summary, and exit with status 1.
\item[-r] The record to create, \textit{RECORD}: its header is \textit{RECORD}\file{.hea}
  (default: the WAV file name without \file{.wav}).
\item[-U] The physical units (required unless the gain is 0).
\end{params}
For example, for the WAV file written by the example of \cmd{mit2wav} below:
\begin{lstlisting}[style=text]
$ wav2mit -i 100s.wav -r w2 -G 200 -U mV
$ cat w2.hea
w2 2 360 21600
100s.wav 16+44 200/mV 16 0 0 0 0
100s.wav 16+44 200/mV 16 0 0 0 0
\end{lstlisting}
\returns{Exit status 0.}
\raises{Exit status 1 for an error in the arguments or a missing \opt{-G} or \opt{-U};
2 if the WAV file cannot be opened; 3 if it is not a supported WAV file (not PCM, more
than 16 bits per sample, inconsistent chunk lengths); 4 if the header cannot be written.}
\textit{Differences from the C tool} (\cref{chap:behavior}): the new options \opt{-G}
(required) and \opt{-U} (required unless the gain is 0); the C tool wrote 12.5\,adu/mV for
8-bit and 6400\,adu/mV for 16-bit files, in millivolts. The \opt{-i} file is read from the
input directory (C: from the current directory).

% ------------------------------------------------------------------------------------------
\apientry{mit2wav}{Convert a record to a WAV file}
\begin{lstlisting}[style=signature,language={}]
usage: mit2wav -o FILE.wav -r RECORD [ OPTIONS ... ]
\end{lstlisting}
Writes the signals of a record to a WAV file (16-bit PCM, one channel per signal, at the
sampling frequency of the record), the samples shifted left to 16 bits (the WAV header gives
the largest ADC resolution as the number of bits per sample). The samples are written
through the library, and the WAV header is then put in front of them, in the same file.
With \opt{-n}, it also writes a header that describes the WAV file as a record, with the
gains and baselines adjusted to the shift.
\begin{params}
\item[-o] The WAV file to write, \textit{FILE.wav} (required; the name must end in
  \file{.wav}), created in the output directory (\cref{sec:tools-files}).
\item[-r] The input record, \textit{RECORD} (required).
\item[-h] Print the usage summary, and exit with status 1.
\item[-n] Write a header for the WAV file, as the record \textit{NEWREC} that follows (it
  must differ from \textit{RECORD}).
\end{params}
\begin{lstlisting}[style=text]
$ mit2wav -r 100s -o 100s.wav -n wavrec
$ cat wavrec.hea
wavrec 2 360 21600
100s.wav 16+44 6400 16 0 -928 -31712 0 MLII
100s.wav 16+44 6400 16 0 -416 4288 0 V5
\end{lstlisting}
\returns{Exit status 0.}
\raises{Exit status 1 for an error in the arguments; 2 if the record cannot be opened; 3 if
its signals cannot be read, it gives no sampling frequency, or the header of \opt{-n}
cannot be written; 4 if the output file cannot be opened.}
\textit{Differences from the C tool} (\cref{chap:behavior}): the record must give its
sampling frequency (C: 250\,Hz assumed); a signal of more than 16 bits is shifted right to
16 bits and gets a gain multiplier of 0 (C: undefined shift, samples written as 0); the WAV
header is put in front of the samples in the file that the library created, in the output
directory (C: in the file of that name in the current directory).

% ------------------------------------------------------------------------------------------
\apientry{wfdb2mat}{Convert a record to a MATLAB/Octave file}
\begin{lstlisting}[style=signature,language={}]
usage: wfdb2mat -r RECORD [OPTIONS ...]
\end{lstlisting}
Writes the samples of a record (in ADC units, less the ADC zero) as a MATLAB/Octave
\file{.mat} file (MATLAB 5 format) holding a matrix \code{val} of one row per signal, and a
header that describes the \file{.mat} file as a record: format 80, 16 or 32 (the smallest
that holds the ADC resolutions of the signals, or their storage formats where the resolution
is not specified) after a 192-byte prolog. Both are named after the record
with an \code{m} appended: \file{100sm.mat} and \file{100sm.hea} for \code{100s} (the help
text, that of the C tool, says \textit{RECORD}\file{.mat}). A summary of the conversion is
printed on the standard output.
\begin{params}
\item[-r] The input record, \textit{RECORD} (required).
\item[-f] Begin at the time that follows.
\item[-h] Print the usage summary, and exit with status 0.
\item[-H] Read multifrequency records in high-resolution mode.
\item[-l] Stop after the time interval that follows (as in \cmd{rdsamp}).
\item[-s] Convert only the signals that follow, by number or name:
  \opt{-s} \textit{SIGNAL} \code{[}\textit{SIGNAL} \code{...]}.
\item[-S] Begin at the first valid sample of the signal that follows, at or after the
  \opt{-f} time.
\item[-t] Stop at the time that follows.
\end{params}
\begin{lstlisting}[style=text]
$ wfdb2mat -r 100s -t 2
Source: record 100s
val has 2 rows (signals) and 720 columns (samples/signal)
Duration:     0:02
Sampling frequency: 360 Hz  Sampling interval: 0.002777777778 sec
Row     Signal  Gain    Base    Units
1       MLII    200     0       nd
2       V5      200     0       nd

To convert from raw units to the physical units shown
above, call the 'rdmat.m' function from the wfdb-matlab
toolbox: https://physionet.org/physiotools/matlab/wfdb-app-matlab/
\end{lstlisting}
\returns{Exit status 0.}
\raises{Exit status 1 for an error in the arguments, an empty interval
(\code{wfdb2mat: no output}), an output larger than 2\,GB, or an output file that cannot be
opened; 2 if the record or a signal cannot be opened, or the
record gives no sampling frequency.}
\textit{Differences from the C tool} (\cref{chap:behavior}): the record must give its
sampling frequency (C: 250\,Hz); uncalibrated signals stay uncalibrated in the output
header (C: 200\,adu/mV); unspecified units are written as \code{nd} (C: \code{mV}); a
one-character record name \code{m} gives the output \code{mm} (C wrote over the input
header).

% ------------------------------------------------------------------------------------------
\apientry{a2m}{Convert AHA-format annotations to WFDB format}
\begin{lstlisting}[style=signature,language={}]
usage: a2m -i AHAFILE [ OPTIONS ... ]
\end{lstlisting}
Converts an AHA-format annotation file (\cref{sec:formats-aha-annotations}) to an MIT-format
annotation file. Four kinds of input are read: files in AHA format written by WFDB tools
(type 0, times in sample intervals), short-format and long-format AHA tape files (types 1
and 2, times in milliseconds), and the compressed files of the AHA Database CD-ROMs and
diskettes (type 3, 6 bytes per annotation, times in milliseconds). The annotation times are
shifted by an offset (\opt{-s}).
\begin{params}
\item[-i] The AHA-format input file, \textit{AHAFILE} (required). It is read from the
  input directory (\cref{sec:tools-files}).
\item[-a] The output annotator (default \code{atr}).
\item[-F] The sampling frequency, used (and then required) if the record header does not
  give one (AHA Database: 250).
\item[-h] Print the usage summary, and exit with status 1.
\item[-r] The output record, \textit{RECORD} (default: the input file name without its
  directory and suffix).
\item[-s] Shift the annotations by the time that follows (default: 0 for type 0, 5:0 for
  type 1, 2:30:0 for type 2, and for type 3 5:0 or 2:30:0 for the records \code{x2xx} and
  \code{x3xx}, or \code{x0xx} and \code{x1xx}).
\item[-t] The type of the input, from 0 to 3, as above (default: 3 if the file name ends
  in \file{.ANO} or \file{.ano}, 0 otherwise).
\end{params}
\begin{lstlisting}[style=text]
$ a2m -i 1201.ANO -F 250
init: can't open header for record 1201
$ rdann -r 1201 -a atr -F 250 -t 5:3
init: can't open header for record 1201
    5:01.000    75250     N    0    0    0
    5:01.812    75453     N    0    0    0
    5:02.628    75657     V    0    0    0
\end{lstlisting}
\returns{Exit status 0.}
\raises{Exit status 1 for an error in the arguments or a sampling frequency that is not
specified; 2 if the input or the annotation files cannot be opened.}
\textit{Differences from the C tool} (\cref{chap:behavior}): the new option \opt{-F}
(the C tool assumed 250\,Hz with a header, or 1 sample per second without one), and
annotation times in milliseconds (types 1 to 3) are converted at the record's sampling
frequency (C: 4\,ms per sample, whatever the frequency); the \opt{-i} file is read from the
input directory (C: from the current directory).

% ------------------------------------------------------------------------------------------
\apientry{m2a}{Convert WFDB annotations to AHA format}
\begin{lstlisting}[style=signature,language={}]
usage: m2a -r RECORD -a MIT-ANNOTATOR AHA-ANNOTATOR [ OPTIONS ... ]
\end{lstlisting}
Copies an MIT-format annotation file to an AHA-format annotation file of the same record
(\cref{sec:formats-aha-annotations}), the annotation types mapped by \code{mamap}. With
\opt{-s}, the annotations before the shift are dropped, and the times of the others are
shifted back and converted to milliseconds, as in the AHA tape files.
\begin{params}
\item[-r] The record, \textit{RECORD} (required).
\item[-a] The input (MIT) and output (AHA) annotators, \textit{MIT-ANNOTATOR}
  \textit{AHA-ANNOTATOR} (required).
\item[-F] The sampling frequency, used if the record header gives none; then required
  for \opt{-s}.
\item[-h] Print the usage summary, and exit with status 1.
\item[-s] Shift the annotation times back by the time that follows, and convert them to
  milliseconds.
\end{params}
\begin{lstlisting}[style=text]
$ m2a -r 100s -a atr aha
\end{lstlisting}
\returns{Exit status 0 (the file \file{100s.aha} is written).}
\raises{Exit status 1 for an error in the arguments or a sampling frequency that is needed
and not specified; 2 if an annotation file cannot be opened.}
\textit{Differences from the C tool} (\cref{chap:behavior}): the new option \opt{-F}; the C
tool assumed 250\,Hz.

% ------------------------------------------------------------------------------------------
\apientry{ad2m}{Convert an AHA-format signal file to a WFDB record}
\begin{lstlisting}[style=signature,language={}]
usage: ad2m -i AHAFILE [ OPTIONS ... ]
\end{lstlisting}
Converts the two signals of an AHA Database signal file, in the old tape format (format 16
with an end-of-data marker) or in the compressed format of the CD-ROMs (\opt{-c}), to a
WFDB record in format 212: \textit{RECORD}\file{.dat} and \textit{RECORD}\file{.hea}.
\begin{params}
\item[-i] The AHA-format input file, \textit{AHAFILE} (required). It is read from the
  input directory (\cref{sec:tools-files}).
\item[-c] The input is compressed (the default if the name ends in \file{.CMP} or
  \file{.cmp}; otherwise the old tape format is assumed).
\item[-f] Begin at the time that follows (default 0, the beginning of the input file).
\item[-F] The sampling frequency in Hz (required; AHA Database: 250).
\item[-G] The gain in adu per physical unit (required; AHA Database: 400; 0 declares the
  signals uncalibrated).
\item[-h] Print the usage summary, and exit with status 1.
\item[-r] The output record, \textit{RECORD} (default: the input file name without its
  directory and suffix).
\item[-R] The ADC resolution in bits, from 1 to 32 (required; AHA Database: 12).
\item[-t] Stop at the time that follows (default: 35 minutes after the beginning for the
  records \code{x2xx} and \code{x3xx}, 3 hours for \code{x0xx} and \code{x1xx}, otherwise
  the end of the input, with a warning).
\item[-U] The physical units (required unless the gain is 0; AHA: \code{mV}).
\end{params}
\begin{lstlisting}[style=text]
$ ad2m -i 1201.CMP -F 250 -G 400 -R 12 -U mV
$ cat 1201.hea
1201 2 250 800
1201.dat 212 400/mV 12 0 0 32104 0
1201.dat 212 400/mV 12 0 90 -31586 0
\end{lstlisting}
\returns{Exit status 0.}
\raises{Exit status 1 for an error in the arguments, a missing \opt{-F}, \opt{-G},
\opt{-R} or \opt{-U}, or an end time not after the start time; 2 if the input cannot be
read; 3 if the output signal file cannot be opened.}
\textit{Differences from the C tool} (\cref{chap:behavior}): the new options \opt{-F},
\opt{-G}, \opt{-R} and \opt{-U}, all required (\opt{-U} unless the gain is 0); the C tool
assumed the parameters of the AHA Database (250\,Hz, 400\,adu/mV, 12 bits, millivolts).
The \opt{-f} and \opt{-t} times are converted at the \opt{-F} frequency. The \opt{-i} file
is read from the input directory (C: from the current directory).

% ------------------------------------------------------------------------------------------
\apientry{md2a}{Convert a record to an AHA-format signal file}
\begin{lstlisting}[style=signature,language={}]
usage: md2a -o AHAFILE -r RECORD [ OPTIONS ... ]
\end{lstlisting}
Writes the signals of a record to an AHA-format signal file: format 16, the signals
interleaved, followed by an end-of-data marker. With \opt{-n}, it also writes a header that
describes the output file as a record.
\begin{params}
\item[-o] The output file, \textit{AHAFILE} (required).
\item[-r] The input record, \textit{RECORD} (required).
\item[-h] Print the usage summary, and exit with status 1.
\item[-n] Write a header for the output file, as the record \textit{NEWREC} that follows.
\end{params}
\begin{lstlisting}[style=text]
$ md2a -r 100s -o 100s.ah -n ahrec
$ cat ahrec.hea
ahrec 2 360 21600
100s.ah 16 200 11 1024 995 21537 0
100s.ah 16 200 11 1024 1011 -3962 0
\end{lstlisting}
\returns{Exit status 0.}
\raises{Exit status 1 for an error in the arguments; 2 if the record cannot be opened; 3 if
the header of \opt{-n} cannot be created; 4 if the output file cannot be opened.}
\textit{Differences from the C tool}: none, apart from \opt{-{}-legacy} (the samples of a
frame beyond those of the record are 0, where the C tool used uninitialized values).

% ------------------------------------------------------------------------------------------
\apientry{ahaecg2mit}{Convert AHA DVD files to WFDB records}
\begin{lstlisting}[style=signature,language={}]
usage: ahaecg2mit [-s] -F FREQ -G GAIN -R BITS -U UNITS RECORD.XXX [RECORD.XXX]...
\end{lstlisting}
Converts the files of the AHA Database DVD, in binary (\file{.ecg}) or text (\file{.txt})
form, to WFDB records: for each input file \textit{RECORD}\file{.XXX}, it writes
\textit{RECORD}\file{.hea}, \textit{RECORD}\file{.dat} (format 212) and
\textit{RECORD}\file{.atr}, and reports \code{wrote ...} on the standard error. Files with
another suffix are ignored, with a message. The options come first, in any order.
\begin{params}
\item[-s] Make short-format (35-minute) records from long-format (2.5-hour) files: the
  first 145 minutes are skipped (located from the \opt{-F} frequency), and the second digit
  of the output record name is incremented by 2.
\item[-F] The sampling frequency in Hz (required; AHA DVD: 250).
\item[-G] The gain in adu per physical unit (required; AHA DVD: 400).
\item[-R] The ADC resolution in bits, from 1 to 32 (required; AHA DVD: 12).
\item[-U] The physical units (required unless the gain is 0; AHA DVD: \code{mV}).
\item[RECORD.XXX] An input file; \textit{RECORD} is a 4-digit AHA Database record name, and
  \textit{XXX} is \code{ecg} or \code{txt}. Several files may be converted in one run. The
  files are read from the input directory (\cref{sec:tools-files}).
\end{params}
\begin{lstlisting}[style=text]
$ ahaecg2mit -F 250 -G 400 -R 12 -U mV 1001.ecg
Rearranging annotations for output annotator atr ...done!
wrote 1001.atr, 1001.dat, and 1001.hea
$ cat 1001.hea
1001 2 250 700
1001.dat 212 400/mV 12 0 0 1213 0 ECG0
1001.dat 212 400/mV 12 0 300 -1026 0 ECG1
\end{lstlisting}
\returns{Exit status 0 (also when an input file is ignored or cannot be opened).}
\raises{Exit status 1 without arguments, with \opt{-h} as the first argument (the usage), or
for a missing or invalid \opt{-F}, \opt{-G}, \opt{-R} or \opt{-U}.}
\textit{Differences from the C tool} (\cref{chap:behavior}): the new options \opt{-F},
\opt{-G}, \opt{-R} and \opt{-U}, required (\opt{-U} unless the gain is 0), which precede the
input files in any order (the C tool had only \opt{-s}, as the first argument, and assumed
250\,Hz, 400\,adu/mV, 12 bits and millivolts); the input files are read from the input
directory (C: from the current directory).

% ------------------------------------------------------------------------------------------
\apientry{makeid}{Write an AHA-format record ID block}
\begin{lstlisting}[style=signature,language={}]
usage: makeid annotator record
\end{lstlisting}
Writes the 512-byte ID block of the AHA tape format for a record (with two signals) and an
annotator to the standard output: the record ID (8 bytes), the number of annotations, the
first and last samples of the annotated segment, the number of 512-byte blocks of samples,
the times of the first and last annotations, and the number of blocks of annotations, as
16- and 32-bit integers. It describes these values on the standard error. The record name
must have at most 8 characters.
\begin{params}
\item[annotator] The annotator.
\item[record] The record.
\end{params}
\begin{lstlisting}[style=text]
$ makeid atr 100s > 100s.id
100s
75 annotations
Annotated segment begins at sample 0
Annotated segment ends at sample 21600
169 blocks of sample data
First annotation is at sample 18
Last annotation is at sample 21423
3 blocks of annotations
\end{lstlisting}
\returns{Exit status 0.}
\raises{Exit status 1 with fewer than two arguments or a record name longer than 8
characters; 2 if the record or the annotator cannot be opened.}
\textit{Differences from the C tool}: none, apart from \opt{-{}-legacy}.

% ------------------------------------------------------------------------------------------
\apientry{readid}{Read an AHA-format record ID block}
\begin{lstlisting}[style=signature,language={}]
usage: readid data-block-size annot-block-size
\end{lstlisting}
Reads an ID block (as written by \cmd{makeid}) from the standard input, and prints the
record ID, the number of data blocks and the number of annotation blocks, recounted in
blocks of the given sizes (in bytes).
\begin{params}
\item[data-block-size] The size of a block of samples, in bytes (at least 1).
\item[annot-block-size] The size of a block of annotations, in bytes (at least 1).
\end{params}
\begin{lstlisting}[style=text]
$ readid 512 16 < 100s.id
100s 169 96
\end{lstlisting}
\returns{Exit status 0.}
\raises{Exit status 1 with fewer than two arguments.}
\textit{Differences from the C tool}: none, apart from \opt{-{}-legacy} (a record ID of
more than 8 characters, which overflows a buffer in the C tool, is printed whole).

% ------------------------------------------------------------------------------------------
\apientry{revise}{Convert an obsolete-format header}
\begin{lstlisting}[style=signature,language={}]
usage: revise record
\end{lstlisting}
Reads the obsolete-format header of a record, \textit{record}\file{.header} (one value per
line, as written by early versions of the WFDB library), and writes a header in the current
format, \textit{record}\file{.hea}.
\begin{params}
\item[record] The record (a \code{header.} prefix is removed).
\end{params}
\returns{Exit status 0.}
\raises{Exit status 1 without an argument; 2 if the old header cannot be read.}
\textit{Differences from the C tool} (\cref{chap:behavior}): a missing ADC resolution stays
unspecified (C: 12 bits).

% ------------------------------------------------------------------------------------------
\apientry{parsescp}{Read and convert an SCP-ECG file}
\begin{lstlisting}[style=signature,language={}]
usage: parsescp -o RECORD [OPTIONS ...] <SCPFILE
\end{lstlisting}
Reads an SCP-ECG file from the standard input, checks its CRCs, decodes its sections
(patient and acquisition data, Huffman tables, leads, QRS locations, reference beats,
rhythm data, measurements, diagnoses) and the ECG, derives the missing limb leads, and
writes: \textit{RECORD}\file{.des}, a description (the patient data, acquisition data and
measurements), \textit{RECORD}\file{.key}, the patient's name and medical record number,
and \textit{RECORD}\file{.ecg}, the 12-lead ECG in a binary layout of its own; with
\opt{-w}, a WFDB record of 12 signals, \textit{RECORD}\file{.hea} and
\textit{RECORD}\file{.dat} (format 16, 200\,adu/mV); with \opt{-v}, the samples as text,
\textit{RECORD}\file{.txt}, and an analysis of the file on the standard output. All these
files are created in the output directory (\cref{sec:tools-files}).
\begin{params}
\item[-o] The record name, \textit{RECORD}, of at most 11 characters (default \code{ecg}).
\item[-a] Anonymize the input and copy it to the standard output (the CRC recomputed);
  no other output.
\item[-f] Parse a file with CRC errors.
\item[-h] Print the usage summary, and exit with status 1.
\item[-l] Low-pass filter (smooth) the output waveforms.
\item[-v] Print a (very) detailed analysis of the input on the standard output, and write
  \textit{RECORD}\file{.txt}; given twice, the text file has column headings.
\item[-w] Write a WFDB-format record.
\item[-z] Make the signals zero-mean.
\item[-b] (Debugging) Show the baseline (residuals), without the reference beats.
\item[-s] (Debugging) Shift the reference beats by the number of samples that follows
  before adding them.
\item[-S] (Debugging) Skip the parsing of the section whose number follows (\code{-} for
  all sections).
\item[-t] (Debugging) Show the reference beats only, without the baseline.
\item[-x] (Debugging) Show hex dumps (implies \opt{-v}).
\end{params}
\begin{lstlisting}[style=text]
$ parsescp -o scp -w < example.scp
$ head -3 scp.hea
scp 12 500 5000     0:22
scp.dat 16 200/mV 16 0 -1 -2500 0 I
scp.dat 16 200/mV 16 0 -3 -2034 0 II
\end{lstlisting}
\begin{phibox}
The \file{.key} file holds the patient's name and medical record number, the \file{.des}
file and the header the age and sex, and \opt{-v} prints all the patient data of the file;
\opt{-a} writes a copy of the input with these fields removed.
\end{phibox}
\returns{Exit status 0.}
\raises{Exit status 1 for an error in the arguments, an input shorter than 6 bytes, a CRC
error (without \opt{-f}), or a reference-beat subtraction zone larger than the template; 2 if the declared length of the
input cannot be allocated; 3 if the input is shorter than its declared length.}
\textit{Differences from the C tool} (\cref{chap:behavior}): the record is written with the
sampling frequency given in the SCP file, and is not written if the file gives none
(\code{parsescp: sampling frequency not specified; record R not written}); the C tool
always wrote 500\,Hz. The \file{.des}, \file{.ecg}, \file{.key} and \file{.txt} files are
created in the output directory, as the record is (C: in the current directory). Signals and
reference beats that cannot be allocated, and samples
beyond the end of a signal shorter than lead 0, are read as absent or as zeros (undefined
behavior in the C tool).
